% Part: normal-modal-logic
% Chapter: syntax-and-semantics
% Section: substitution

\documentclass[../../../include/open-logic-section]{subfiles}

\begin{document}

\olfileid{nml}{syn}{sub}

\olsection{એકસાથે પ્રતિસ્થાપન}

!!a{formula}~$!A$માં !!a{propositional variable}નાં બધાં આવર્તનોને
કોઈ બીજા !!{formula}થી બદલીને મળતું પરિણામ એ~$!A$નો
\emph{દૃષ્ટાંત} છે. પ્રામાણ્ય અને !!{derivability} બંનેની ચર્ચામાં
!!{formula}sના આવા દૃષ્ટાંતોનો વારંવાર ઉલ્લેખ આવશે. તેથી આ ખ્યાલની
ચોક્કસ વ્યાખ્યા આપવી ઉપયોગી છે.

\begin{defn}\ollabel{def:subst-inst}
  ધારો કે $!A$ એ મોડલ !!{formula} છે જેમાંનાં બધાં !!{propositional
    variable}s $p_1$, \dots, $p_n$ પૈકીનાં છે, અને $!D_1$, \dots,
  $!D_n$ પણ મોડલ !!{formula}s છે. ત્યારે
  $\SSubst{!A}{\subst{!D_1}{p_1}, \dots, \subst{!D_n}{p_n}}$ને
  $!A$માં દરેક $p_i$ માટે એકસાથે $!D_i$નું પ્રતિસ્થાપન
  કરવાથી મળતું પરિણામ વ્યાખ્યાયિત કરીએ છીએ. ચોક્કસ રીતે,
  આ વ્યાખ્યા~$!A$ પર અનુમાન વડે અપાય છે:
  \begin{enumerate}
    \tagitem{prvFalse}{\indcase{!A}{\lfalse}{%
        $\SSubst{\indfrm}{\subst{!D_1}{p_1}, \dots,
          \subst{!D_n}{p_n}}$ એ $\lfalse$ છે.}}{}
    \tagitem{prvTrue}{\indcase{!A}{\ltrue}{%
        $\SSubst{\indfrm}{\subst{!D_1}{p_1}, \dots,
          \subst{!D_n}{p_n}}$ એ $\ltrue$ છે.}}{}
    \item \indcase{!A}{q}{$\SSubst{\indfrm}{\subst{!D_1}{p_1}, \dots,
        \subst{!D_n}{p_n}}$ એ $q$ છે, જો $i
      = 1$, \dots,~$n$ માટે $q \not\ident p_i$ હોય.}
    \item \indcase{!A}{p_i}{$\SSubst{\indfrm}{\subst{!D_1}{p_1},
        \dots, \subst{!D_n}{p_n}}$ એ $!D_i$ છે.}
    \tagitem{prvNot}{\indcase{!A}{\lnot !B}{%
        $\SSubst{\indfrm}{\subst{!D_1}{p_1}, \dots, \subst{!D_n}{p_n}}$ એ
        $\lnot \SSubst{!B}{\subst{!D_1}{p_1}, \dots, \subst{!D_n}{p_n}}$ છે.}}{}
    \tagitem{prvAnd}{\indcase{!A}{(!B \land
        !C)}{$\SSubst{\indfrm}{\subst{!D_1}{p_1}, \dots,
          \subst{!D_n}{p_n}}$ એ
        \[(\SSubst{!B}{\subst{!D_1}{p_1}, \dots, \subst{!D_n}{p_n}} \land
          \SSubst{!C}{\subst{!D_1}{p_1}, \dots, \subst{!D_n}{p_n}}).\] છે.}}{}
    \tagitem{prvOr}{\indcase{!A}{(!B \lor
          !C)}{$\SSubst{\indfrm}{\subst{!D_1}{p_1}, \dots,
          \subst{!D_n}{p_n}}$ એ
        \[(\SSubst{!B}{\subst{!D_1}{p_1}, \dots, \subst{!D_n}{p_n}} \lor
          \SSubst{!C}{\subst{!D_1}{p_1}, \dots, \subst{!D_n}{p_n}}).\] છે.}}{}
    \tagitem{prvIf}{\indcase{!A}{(!B \lif
          !C)}{$\SSubst{\indfrm}{\subst{!D_1}{p_1}, \dots,
          \subst{!D_n}{p_n}}$ એ
        \[(\SSubst{!B}{\subst{!D_1}{p_1}, \dots, \subst{!D_n}{p_n}} \lif
          \SSubst{!C}{\subst{!D_1}{p_1}, \dots, \subst{!D_n}{p_n}}).\] છે.}}{}
    \tagitem{prvIff}{\indcase{!A}{(!B \liff
          !C)}{$\SSubst{\indfrm}{\subst{!D_1}{p_1}, \dots,
          \subst{!D_n}{p_n}}$ એ
        \[(\SSubst{!B}{\subst{!D_1}{p_1}, \dots, \subst{!D_n}{p_n}} \liff
          \SSubst{!C}{\subst{!D_1}{p_1}, \dots, \subst{!D_n}{p_n}}).\] છે.}}{}
    \tagitem{prvBox}{\indcase{!A}{\Box !B}{%
        $\SSubst{\indfrm}{\subst{!D_1}{p_1}, \dots, \subst{!D_n}{p_n}}$ એ
        $\Box
        \SSubst{!B}{\subst{!D_1}{p_1}, \dots, \subst{!D_n}{p_n}}.$ છે}}{}
    \tagitem{prvDiamond}{\indcase{!A}{\Diamond !B}{%
        $\SSubst{\indfrm}{\subst{!D_1}{p_1}, \dots, \subst{!D_n}{p_n}}$ એ
        $\Diamond
        \SSubst{!B}{\subst{!D_1}{p_1}, \dots, \subst{!D_n}{p_n}}.$ છે}}{}
  \end{enumerate}
  !!{formula} $\SSubst{!A}{\subst{!D_1}{p_1}, \dots,
    \subst{!D_n}{p_n}}$ને $!A$નો \emph{પ્રતિસ્થાપન દૃષ્ટાંત}
  કહે છે.
\end{defn}
\sourcecorrection{OLMD-002}{સ્થિર અંગ્રેજી મૂળ
દ્વિમુખી પ્રેરણના આગમન-કિસ્સાને ભૂલથી
\texttt{prvIf} ટૅગ આપે છે, જ્યારે તેની સૂત્રરચના
$!B \liff !C$ છે. મૂળ ભાષાની રચના-વ્યાખ્યા
સાથે સુસંગત \texttt{prvIff} ટૅગ અહીં રાખ્યો છે.}
\sourcecorrection{OLMD-003}{સ્થિર અંગ્રેજી મૂળ
$\Box !B$ના આગમન-કિસ્સાને અશરતી \texttt{item}
રૂપે રાખે છે, જ્યારે મૂળ ભાષામાં $\Box$નો પ્રાથમિક
કારક \texttt{prvBox} ટૅગથી પસંદ થાય છે. તેની સાથે
સુસંગત રીતે આ કિસ્સો \texttt{prvBox}થી શરતિત કર્યો છે.}

\begin{ex}
  માનો કે $!A$ એ $p_1 \lif \Box(p_1 \land p_2)$ છે,
  $!D_1$ એ $\Diamond(p_2 \lif p_3)$ છે અને $!D_2$
  એ $\lnot\Box p_1$ છે. ત્યારે
  $\SSubst{!A}{\subst{!D_1}{p_1}, \subst{!D_2}{p_2}}$ આ છે:
  \begin{align*}
    \Diamond(p_2 \lif p_3) &
    \lif \Box(\Diamond(p_2 \lif p_3) \land \lnot\Box p_1)
    \intertext{જ્યારે $\SSubst{!A}{\subst{!D_2}{p_1}, \subst{!D_1}{p_2}}$ આ છે:}
    \lnot\Box p_1 &
    \lif \Box(\lnot\Box p_1 \land \Diamond(p_2 \lif p_3))
    \intertext{નોંધો કે એકસાથે પ્રતિસ્થાપન સામાન્ય રીતે
      ક્રમિક પ્રતિસ્થાપન જેવું નથી. દાખલા તરીકે, ઉપરના
      $\SSubst{!A}{\subst{!D_1}{p_1}, \subst{!D_2}{p_2}}$ને
      $\Subst{(\Subst{!A}{!D_1}{p_1})}{!D_2}{p_2}$ સાથે સરખાવો; તે આ છે:}
    \Diamond(p_2 \lif p_3) & \Subst{\lif \Box(\Diamond(p_2 \lif p_3) \land p_2)}{\lnot\Box p_1}{p_2}, \text{ અર્થાત્,}\\
    \Diamond(\lnot\Box p_1 \lif p_3) &
    \lif \Box(\Diamond(\lnot\Box p_1
    \lif p_3) \land \lnot\Box p_1)\\
    \intertext{અને $\Subst{(\Subst{!A}{!D_2}{p_2})}{!D_1}{p_1}$ સાથે પણ સરખાવો:}
    p_1 & \lif \Subst{\Box(p_1 \land \lnot\Box p_1)}{\Diamond(p_2 \lif p_3)}{p_1}, \text{ અર્થાત્,}\\
    \Diamond(p_2 \lif p_3) &
    \lif \Box(\Diamond(p_2
    \lif p_3) \land \lnot\Box\Diamond(p_2 \lif p_3)).
  \end{align*}
\end{ex}

\end{document}
